% =========================================================
% Proof: Associativity of Symmetric Difference
% Source: volume-i/algebras-of-sets/notes/...
% Mode: standard
% =========================================================

\subsection*{Associativity of Symmetric Difference}
\label{prf:symmdiff-assoc}

\begin{remark*}[Return]
\hyperref[prop:symmdiff-assoc]{$\leftarrow$ Back to Proposition (Associativity of Symmetric Difference) in Notes}
\end{remark*}

\begin{proof}
% TODO: proof to be filled.
% Strategy: Parity argument: $x\in A\triangle(B\triangle C)$ iff $x$ belongs to an odd number of $A$, $B$, $C$.
% This parity characterisation is symmetric in how the sets are grouped, giving associativity.
\end{proof}

\begin{remark*}[Proof shape]
% TODO

\FlashProof{Key lemma: $x\in A\triangle B$ iff $x$ is in exactly one of $A$, $B$. Use this to show $x\in A\triangle(B\triangle C)$ iff $x$ is in an odd number of $\{A,B,C\}$. The same parity condition holds for $x\in(A\triangle B)\triangle C$. Since parity is independent of grouping, the two sides agree.}
\end{remark*}

\begin{remark*}[Dependencies]
\begin{itemize}
  \item \hyperref[def:set-operations]{Definition of symmetric difference}
  \item \hyperref[prop:symmdiff-comm]{Commutativity of Symmetric Difference}
\end{itemize}

\FlashUses{def:set-operations, prop:symmdiff-comm}
\end{remark*}
